% =============================================================
% conventions.tex — Conventions de notation
% Légende unique des familles d'identifiants employées dans le corps
% du mémoire, et convention de désignation des composants.
% Voir PLAN_REFONTE_PFE.md §11.6 et §11.7.
% =============================================================
\chapter*{Conventions de notation}
\addcontentsline{toc}{chapter}{Conventions de notation}
\thispagestyle{plain}

\noindent\textbf{Désignation des composants.} Les composants du socle sont désignés par leur
\emph{capacité} --- ce qu'ils font --- plutôt que par leur nom commercial, afin que
l'architecture reste lisible indépendamment du fournisseur. Le produit qui réalise chaque
capacité est indiqué entre parenthèses à sa première occurrence, et la correspondance complète
figure au tableau d'inventaire technologique de la \cref{sec:trente-trois-briques}.

\medskip
\noindent\textbf{Familles d'identifiants.} Le mémoire emploie des codes préfixés pour rendre la
traçabilité vérifiable d'un chapitre à l'autre : une menace, l'exigence qui la traite, le contrôle
qui la met en \oe uvre et le test qui la vérifie portent chacun un identifiant stable. Chaque
famille est \emph{établie} une seule fois, dans le chapitre porté en bandeau du tableau ci-dessous,
et le reste du mémoire n'y renvoie plus que par le code. Les séries sont complètes : aucun numéro
n'est manquant à l'intérieur d'une série.

\begingroup
\setlength{\LTpre}{6pt}
\setlength{\LTpost}{6pt}
\begin{longtable}{|P{1.7cm}|P{10.9cm}|P{2.7cm}|}
\hline
\textbf{Préfixe} & \textbf{Signification} & \textbf{Série} \\
\hline
\endhead
\multicolumn{3}{|l|}{\textbf{Chapitre~\ref{ch:cadre-existant}} --- contexte, existant et problématique} \\ \hline
C & Carence structurante de l'existant & C1--C5 \\ \hline
Q & Question dérivée de la problématique & Q1--Q6 \\ \hline
O & Objectif du projet & O1--O6 \\ \hline
PH & Phase du projet & PH0--PH7 \\ \hline
\multicolumn{3}{|l|}{\textbf{Chapitre~\ref{ch:etat-art}} --- état de l'art et choix technologiques} \\ \hline
V & Verrou identifié dans l'état de l'art & V1--V4 \\ \hline
\multicolumn{3}{|l|}{\textbf{Chapitre~\ref{ch:besoins-menaces}} --- besoins et modélisation des menaces} \\ \hline
A & Acteur du modèle, légitime ou hostile & A1--A11 \\ \hline
BF & Besoin fonctionnel & BF1--BF13 \\ \hline
BNF & Besoin non fonctionnel & BNF1--BNF7 \\ \hline
TB & Frontière de confiance & TB1--TB5 \\ \hline
SO & Scénario opérationnel de menace & SO1--SO24 \\ \hline
EX & Exigence de sécurité & EX1--EX20 \\ \hline
\multicolumn{3}{|l|}{\textbf{Chapitre~\ref{ch:conception}} --- conception de l'architecture} \\ \hline
PR & Principe directeur d'architecture & PR1--PR4 \\ \hline
L & Niveau du modèle en couches & L1--L7 \\ \hline
F & Flux de l'architecture & F1--F7 \\ \hline
D & Décision d'architecture & D00--D16 \\ \hline
É & Écart entre conception et implémentation & É1--É11 \\ \hline
\multicolumn{3}{|l|}{\textbf{Chapitre~\ref{ch:realisation}} --- réalisation et industrialisation} \\ \hline
R & Règle de détection & R1--R7 \\ \hline
\multicolumn{3}{|l|}{\textbf{Chapitre~\ref{ch:validation}} --- validation et mesure} \\ \hline
T & Protocole de la campagne de validation & T1--T20 \\
\hline
\end{longtable}
\endgroup
% Ce tableau est hors numérotation (norme ESPRIT, C26). longtable incrémente le
% compteur « table » même sans \caption : sans ce décrément, la liste des
% tableaux démarrerait au numéro 2.
\addtocounter{table}{-1}

\medskip
\noindent Quatre familles ne servent qu'à l'analyse de risque et sont définies là où elles servent,
en annexe~\ref{ann:modele-menaces} : valeurs métier VM1--VM6, biens supports BS1--BS9, sources de
risque SR1--SR6 et scénarios stratégiques SS1--SS7. Deux échelles de cotation y emploient les
préfixes G et V --- gravité G1--G4, vraisemblance V1--V4 --- sans rapport avec les verrous V1--V4
du tableau ci-dessus : le contexte est toujours celui d'une case de la cartographie du risque. Les
sept jalons M1--M7 du planning sont, eux, définis avec la figure qui les porte, au
chapitre~\ref{ch:realisation}.

% --- Compacité de la table des matières (norme ESPRIT §H : 2 pages) ---
% report.cls insère 1,0 em de blanc avant chaque entrée de niveau chapitre
% (\l@chapter, l. 615). Le mémoire en compte vingt-trois : ce seul interligne
% coûte une page de liminaires. Il est ramené à 0,1 em ; les entrées de section
% ne sont pas touchées. Emplacement provisoire : ce réglage appartient à
% config/formatting.tex, hors du périmètre d'écriture de cette équipe ; il est
% posé ici parce que conventions.tex est le dernier fichier chargé avant
% \tableofcontents. À déplacer à la prochaine passe.
\makeatletter
\patchcmd{\l@chapter}{\vskip 1.0em \@plus\p@}{\vskip 0.1em \@plus\p@}{}%
  {\typeout{ATTENTION : patch de compacite de la table des matieres NON applique}}
\makeatother

\clearpage
